\unitopen{47}
\begin{candos}
\item czytać artykuł prasowy i techniczny tekst i znaleźć główną tezę i szczegóły w limicie czasu;
\item słuchać wiadomości i podcastów w standardzie oraz fragmentów w dialekcie i rozumieć sedno;
\item stosować strategie egzaminacyjne (słowa kluczowe, eliminowanie, notowanie);
\item prowadzić dziennik słuchania i czytania i mierzyć postęp.
\end{candos}

\section{Słownictwo}
\vocabtable{47}

\section{Czytanie: trzy przebiegi}
\textbf{Przebieg 1 (30 sekund): struktura.} \de{Überschrift}, pierwszy i ostatni akapit, liczby i nazwy własne. \textbf{Przebieg 2 (dokładny):} akapit po akapicie
zapisz słowo kluczowe i jedno zdanie; rozpoznaj odniesienia (\de{dies}, \de{dabei}, \de{dadurch}). \textbf{Przebieg 3 (sprawdzenie):} wróć do
pytań i znajdź \emph{dowód} w tekście; nie wybieraj odpowiedzi „bo brzmi dobrze”.

\section{Słuchanie: strategie}
Przed nagraniem przeczytaj zadania i przewidź słownictwo. Notuj liczby, imiona, negacje (\de{nicht}, \de{kein}, \de{nie}) i spójniki przeciwstawne
(\de{aber}, \de{jedoch}, \de{allerdings}): często przenoszą sens. Nie zatrzymuj się na nieznanym słowie. Dwukrotne słuchanie: najpierw sedno, potem szczegóły.

\begin{center}\small
\begin{tabular}{@{}p{0.2\linewidth}p{0.72\linewidth}@{}}
\toprule
\textsf{\textbf{Źródło}} & \textsf{\textbf{Zastosowanie}}\\
\midrule
\de{Tagesschau in 100 Sekunden} & krótkie wiadomości w standardzie, DE\\
\de{DW Langsam gesprochene Nachrichten} & wolne wiadomości, B1--B2\\
\de{SRF News}, \de{SRF Meteo} & szwajcarski standard i dialekt, CH\\
\de{Podcasts} z IT (po niemiecku) & słownictwo zawodowe\\
\de{NZZ}, \de{Tages-Anzeiger}, \de{Die Zeit} & czytanie, B2\\
\bottomrule
\end{tabular}
\end{center}
Sprawdź aktualność źródeł; nazwy audycji i oferta zmieniają się.

\begin{zasada}[Dziennik]
Co tydzień: 3 teksty przeczytane i 3 nagrania wysłuchane. Zapisz źródło, temat, pięć nowych słów, jedno zdanie streszczenia i ocenę
zrozumienia w procentach. Po czterech tygodniach zobaczysz trend; to jedyny wiarygodny pomiar, bo liczba godzin sama niczego nie dowodzi.
\end{zasada}

\begin{falle}
Typowy błąd słuchania: tłumaczenie w głowie słowo po słowie, więc tracisz kolejne zdania. Drugi: wybieranie odpowiedzi, która zawiera słowo
z nagrania (pułapka dystraktora). Trzeci: czytanie ze słownikiem w pierwszym przebiegu. Czwarty: słuchanie wyłącznie ze znanymi napisami
(wyłącz je co tydzień dłużej).
\end{falle}

\begin{szwajcaria}
W Szwajcarii słuchaj \de{SRF} (wiadomości w standardzie, programy w dialekcie), by oswoić ucho. Z dialektu rozumiej najpierw
pytania i liczby (\de{wie vil}, \de{um wie viel Uhr}). Szwajcarzy mówią szybciej, niż się wydaje; nie jest to wada Twojego słuchu, tylko trening, którego jeszcze nie masz.
\end{szwajcaria}

\begin{praca}
Posłuchaj nagrania spotkania technicznego i zrób notatkę w czterech punktach: kto, co zdecydowano, jakie terminy, jakie pytania pozostały.
Przeczytaj dokumentację w języku niemieckim (np.\ strona producenta chmury) i napisz trzy zdania streszczenia.
\end{praca}

\section{Dialogi}
\textbf{Dialog 1 --- dyskusja o artykule}
\begin{dialog}
\turn{A}{Hast du den Artikel über die neue Regulierung gelesen?}{Czytałeś artykuł o nowej regulacji?}
\turn{B}{Ja, die Hauptaussage ist, dass Unternehmen ihre Daten künftig besser schützen müssen. Das Detail zur Frist war mir aber unklar.}{Tak, główna teza jest taka, że firmy będą musiały lepiej chronić dane. Szczegół dotyczący terminu był dla mnie niejasny.}
\turn{A}{Ich habe die Frist herausgehört: Sie läuft bis Ende des Jahres.}{Wychwyciłem termin ze słuchu: upływa z końcem roku.}
\turn{B}{Danke. Dann kann ich die Notiz ergänzen.}{Dzięki. Wtedy mogę uzupełnić notatkę.}
\end{dialog}

\section{Ćwiczenia}
\begin{cwiczenia}
\exitem{Podaj trzy kroki czytania artykułu.}{(1) struktura i nagłówek, (2) dokładne czytanie akapit po akapicie, (3) sprawdzenie odpowiedzi w tekście}
\exitem{Co robisz przed nagraniem?}{Czytam zadania i przewiduję słownictwo.}
\exitem{Które słowa w nagraniu często przenoszą sens?}{negacje i spójniki przeciwstawne: \de{nicht, kein, aber, jedoch, allerdings}}
\exitem{Co to jest pułapka dystraktora?}{odpowiedź zawiera słowo z nagrania, ale jest błędna}
\exitem{Co zawiera wpis w dzienniku słuchania?}{źródło, temat, pięć nowych słów, streszczenie, ocena zrozumienia}
\exitem{Przetłumacz: „Główna teza artykułu jest taka, że \dots”}{\de{Die Hauptaussage des Artikels ist, dass \dots}}
\exitem{Przetłumacz: „Chodzi o \dots”}{\de{Es geht um \dots}}
\exitem{Przetłumacz: „Tekst sugeruje, że \dots”}{\de{Der Text legt nahe, dass \dots}}
\exitem{Zapytaj: „Które stwierdzenie pasuje do tekstu?”}{\de{Welche Aussage passt zum Text?}}
\exitem{Co robisz, gdy nie znasz słowa w nagraniu?}{Nie zatrzymuję się, słucham dalej i wnioskuję z kontekstu.}
\end{cwiczenia}

\weekend{Zrób pełny egzamin próbny z czytania i słuchania z arkusza modelowego B2 (Goethe, telc lub ÖSD) w warunkach egzaminu. Zapisz wynik w procentach,
błędy sklasyfikuj (słownictwo, negacja, dystraktor, czas) i zaplanuj powtórki.}
